\begin{abstract}
Vision-language models (VLMs) have advanced rapidly, yet their ability to perform grounded multimodal legal reasoning remains underexplored. In traffic-law applications, this capability requires models to identify violations from visual evidence, link those observations to applicable legal provisions, and determine the corresponding legal consequences. However, existing traffic-oriented benchmarks typically evaluate visual scene understanding and traffic-law knowledge in isolation, leaving it unclear whether VLMs can translate visual evidence into legally valid decisions. To address this gap, we introduce \textsc{ViTALBench}, a multiple-choice visual question answering benchmark containing 4,536 questions, grounded in current Vietnamese traffic regulations.
\textsc{ViTALBench} contains two top-level tiers: detection and law. The law tier is divided into violation identification, regulation application, and multi-step legal reasoning, evaluating visual scene understanding and statute-dependent reasoning grounded in external legal knowledge. Across ten VLMs from three model families, we find a consistent performance gap between visual perception and legally grounded reasoning, with every model performing substantially worse on questions that require statutory grounding. Further analysis shows that the main bottlenecks are not visual perception alone, but limited statutory knowledge and difficulty integrating legal text with visual evidence to produce legal conclusions.
The dataset is available at \url{https://anonymous.4open.science/r/ViTAL_BENCH-BB6C/}.

\keywords{Vision-language models \and Multimodal legal reasoning \and Traffic law \and Visual question answering \and Benchmark.}
\end{abstract}